\subsection{Milestone 2 -- Identificazione di modelli accurati}

La Milestone 2 estende il dataset prodotto dalla Milestone 1 con una valutazione
comparativa di classificatori per la predizione della bugginess di classi Java in Apache
OpenJPA. Per ogni combinazione di classificatore, strategia di feature selection e
strategia di balancing viene calcolata una serie di metriche di performance tramite
10-times 10-fold cross-validation. I risultati sono scritti in un file CSV dedicato.
La Milestone 2 è eseguibile in modo completamente \textbf{indipendente} dalla Milestone
1: legge il CSV già prodotto da disco senza rieseguire la pipeline.

\subsubsection{Dataset di input}

Il CSV prodotto dalla Milestone 1 contiene una riga per ogni coppia (classe Java,
release), con 19 feature di processo e prodotto e un'etichetta binaria \texttt{isBuggy}
(\texttt{true}/\texttt{false}). Le feature includono metriche di dimensione (LOC,
LOC\_Touched), attività (NR, Nauth, Churn), esperienza degli sviluppatori (EXP, SEXP),
accoppiamento strutturale (Change\_Set\_Size) e code smell (Previous\_Release\_Code\_Smells).

Le colonne \texttt{java\_class\_path} e \texttt{release} sono identificatori, non feature
predittive: includere percorsi di file o nomi di versione nel training farebbe memorizzare
al classificatore etichette associate a stringhe specifiche anziché pattern generalizzabili.
Vengono rimosse prima del caricamento in Weka tramite un filtro \texttt{Remove}.

Il dataset presenta una forte sbilanciamento: circa l'\textbf{8,98\% delle istanze è
buggy} ($\approx$1.378 su $\approx$15.338), rendendo il bilanciamento una fase critica
del preprocessing.

\subsubsection{Classificatori}

Sono stati scelti tre classificatori con approcci radicalmente diversi: uno basato su
ensemble di alberi (RandomForest), uno probabilistico (NaiveBayes) e uno basato su
distanza (IBk). Questa diversità rende il confronto significativo e permette di isolare
quali caratteristiche del modello contano di più per questo problema specifico. L'output
di ciascun classificatore è una distribuzione di probabilità sulle classi, indispensabile
per le metriche AUC e NPofB20 che si basano sul ranking delle istanze per rischio.

\paragraph{RandomForest}

Ensemble di alberi decisionali addestrati ciascuno su un campione bootstrap del training
set, con selezione casuale di un sottoinsieme di feature a ogni split interno. La
predizione finale è la media delle probabilità di tutti gli alberi: questa aggregazione
riduce la varianza tipica di un singolo albero (bagging). L'implementazione usa
\texttt{weka.classifiers.trees.RandomForest} con 100 alberi e $\sqrt{m}$ feature per
split (valori di default Weka). Il numero di thread per la costruzione degli alberi è
configurabile tramite il parametro \texttt{rfSlots}.

\paragraph{NaiveBayes}

Classificatore probabilistico generativo che stima $P(\text{classe} \mid \text{feature})$
tramite il teorema di Bayes, assumendo che tutte le feature siano \textbf{condizionalmente
indipendenti} data la classe. Per feature continue assume una distribuzione gaussiana per
ogni attributo all'interno di ogni classe. L'assunzione di indipendenza è quasi sempre
violata in pratica (LOC\_Touched, LOC\_Added e Churn sono fortemente correlate), ma
NaiveBayes è utile come riferimento e come classificatore interno nei metodi wrapper grazie
alla sua velocità.

\paragraph{IBk (k-Nearest Neighbors)}

Classificatore instance-based (lazy learner) che non costruisce un modello esplicito
durante il training. La predizione di una nuova istanza avviene calcolando la distanza
euclidea da tutte le istanze del training set, trovando i $k$ più vicini ($k=1$, default
Weka) e votando per la classe più frequente. IBk è utile per catturare pattern locali
nello spazio delle feature senza assumere nessuna forma globale della distribuzione, e
fornisce un riferimento non parametrico a confronto con gli approcci parametrici (NaiveBayes)
ed ensemble (RandomForest).

\begin{table}[h!]
\centering
\small
\begin{tabular}{@{}p{0.14\textwidth}p{0.38\textwidth}p{0.38\textwidth}@{}}
\toprule
\textbf{Classificatore} & \textbf{Vantaggi} & \textbf{Svantaggi} \\
\midrule
RandomForest &
Robusto all'overfitting; probabilità ben calibrate (utile per NPofB20); non richiede
normalizzazione; stabile tra run diversi &
Meno interpretabile; il più lento senza parallelismo
($\approx$757\,s per combinazione, rfSlots=1) \\
\addlinespace
NaiveBayes &
Il più veloce ($\approx$6\,s); interpretabile analiticamente;
beneficia dalla feature selection che riduce le ridondanze &
Assunzione di indipendenza quasi sempre violata;
assume distribuzione gaussiana; probabilità spesso mal calibrate \\
\addlinespace
IBk &
Non parametrico; cattura pattern locali complessi;
nessuna fase di training esplicita &
Lento in prediction ($\approx$227\,s); sensibile alla scala delle feature;
con dataset sbilanciato la maggioranza domina i $k$ vicini \\
\bottomrule
\end{tabular}
\caption{Confronto tra i tre classificatori utilizzati. I tempi si riferiscono a
10$\times$10 CV su OpenJPA ($\approx$15.338 istanze) senza balancing.}
\end{table}

\subsubsection{Feature Selection}

La \textbf{feature selection} è il processo di identificazione del sottoinsieme di
attributi più informativi per la predizione, scartando quelli irrilevanti o ridondanti.
Molte delle 19 feature sono fortemente correlate tra loro: \texttt{LOC\_Touched},
\texttt{LOC\_Added} e \texttt{Churn} misurano varianti della stessa grandezza;
\texttt{Max\_LOC\_Added} e \texttt{Avg\_LOC\_Added} sono statistiche derivate dalla
stessa distribuzione. Questa ridondanza ha effetti negativi differenziati per
classificatore: NaiveBayes viola l'assunzione di indipendenza condizionale, IBk soffre
della ``maledizione della dimensionalità'', RandomForest è relativamente robusto grazie
alla selezione casuale delle feature per split.

Il vero vantaggio della feature selection non è necessariamente migliorare le prestazioni
del modello, bensì mantenere prestazioni comparabili con meno feature, ottenendo benefici
concreti: velocità di inference, interpretabilità del modello e minori costi di
manutenzione delle metriche.

Esistono due famiglie di metodi. I metodi \textbf{filter} valutano ogni attributo
indipendentemente dal classificatore che verrà addestrato: sono veloci e general-purpose.
I metodi \textbf{wrapper} valutano sottoinsiemi addestrando il classificatore su ciascun
sottoinsieme candidato: trovano la combinazione ottimale per il classificatore specifico
a costo di tempi molto più elevati.

\paragraph{NONE -- Nessuna feature selection}

Tutti gli attributi vengono passati al classificatore senza alcun preprocessing di
selezione. Fornisce la baseline di confronto: se un metodo di FS non supera NONE in
nessuna configurazione, non ha portato valore aggiunto su questo dataset.

\paragraph{Information Gain (filtro)}

Misura la riduzione di entropia della distribuzione della classe target ottenuta
conoscendo il valore di un attributo:

\[
\text{InfoGain}(C, A) = H(C) - H(C \mid A)
\]

dove l'entropia della classe è $H(C) = -\sum_{c \in C} P(c) \log_2 P(c)$. Per OpenJPA
(9\% buggy): $H(C) \approx 0{,}44$ bit. Un attributo con alto InfoGain riduce molto
l'incertezza sulla bugginess; un attributo irrilevante ha InfoGain $= 0$.

Implementazione: \texttt{InfoGainAttributeEval} + \texttt{Ranker(threshold=0.0)}, che
elimina solo gli attributi con InfoGain esattamente zero. Applicato esclusivamente sul
training fold tramite \texttt{FilteredClassifier}: le label del test non influenzano la
selezione. Overhead computazionale: trascurabile.

\paragraph{Correlazione di Spearman (filtro)}

Misura la correlazione monotona (non necessariamente lineare) tra ogni attributo e la
classe. Il coefficiente $\rho$ è la correlazione di Pearson calcolata sui ranghi delle
osservazioni:

\[
\rho(X, C) = \text{Pearson}(\text{rank}(X),\, \text{rank}(C))
\]

A differenza di Pearson, non assume una relazione lineare ed è robusta agli outlier ---
cruciale per metriche come \texttt{LOC\_Touched} o \texttt{Churn} con distribuzioni
fortemente skewed. Il valore assoluto $|\rho|$ è usato come score perché la direzione
della correlazione non è rilevante per la selezione: una feature fortemente anti-correlata
è ugualmente informativa di una fortemente correlata.

\texttt{weka-stable 3.8.6} include \texttt{CorrelationAttributeEval} (Pearson) ma non un
\texttt{AttributeEvaluator} per Spearman. La classe \texttt{SpearmanAttributeEval} è
un'implementazione originale scritta da zero, basata sulla definizione matematica standard
(Spearman, 1904), con gestione dei pareggi tramite rank medio. Configurazione:
\texttt{SpearmanAttributeEval} + \texttt{Ranker(threshold=0.0)}.

\paragraph{ForwardSearch e BackwardSearch (wrapper)}

I metodi wrapper valutano \textbf{sottoinsiemi} di attributi addestrando il classificatore
su ciascun sottoinsieme candidato e misurando la performance tramite cross-validation
interna. Trovano la combinazione di feature che massimizza la performance del
classificatore specifico, catturando le interazioni tra attributi che i metodi filter
ignorano. Il classificatore interno usato per valutare i subset candidati è NaiveBayes
(il più veloce disponibile).

\textbf{ForwardSearch}: parte con zero feature e aggiunge iterativamente la feature che
produce il maggior miglioramento, fino a quando nessuna aggiunta porta miglioramenti.
Preferibile quando si sospetta che molte feature siano irrilevanti.

\textbf{BackwardSearch}: parte con tutte le feature e rimuove iterativamente quella la
cui rimozione causa il minor degrado, fino a quando nessuna rimozione porta miglioramenti.
Preferibile quando la maggioranza delle feature è utile e solo poche sono irrilevanti.

Entrambi usano \texttt{WrapperSubsetEval(NaiveBayes)} + \texttt{GreedyStepwise},
applicati internamente per ogni fold tramite \texttt{FilteredClassifier}: il loop di
ricerca del subset ottimale opera esclusivamente sui dati del training fold, senza alcuna
visibilità sul test fold.

\medskip
\noindent\textbf{Nota sui tempi:} i metodi wrapper sono \textbf{disabilitati di default}
in \texttt{config.yaml}. Con 19 feature e 10$\times$10 fold CV esterna, ForwardSearch
valuta al massimo $19 + 18 + \ldots + 1 = 190$ subset per fold esterno, ciascuno con
5-fold CV interna su NaiveBayes, per un totale di $\approx$95.000 addestramenti per
combinazione. Il tempo stimato è di 1--3 ore per singola combinazione (24 combinazioni
wrapper totali).

\subsubsection{Bilanciamento delle classi}

Con $\approx$9\% di istanze buggy, un classificatore che predice sempre ``non buggy''
ottiene $\approx$91\% di accuracy senza aver appreso nulla di utile. Il balancing
impedisce al classificatore di imparare questa scorciatoia rendendo le due classi
ugualmente frequenti nel training set, costringendo il modello a imparare i pattern reali
della classe buggy.

Il balancing viene applicato \textbf{esclusivamente sul training fold} tramite
\texttt{FilteredClassifier}: il test fold non viene mai modificato perché deve essere
\textbf{come il mondo reale}, cioè riflettere la distribuzione effettiva del problema
--- quella che il modello incontrerà in produzione ($\approx$9\% buggy, $\approx$91\%
non-buggy).

Applicare balancing o feature selection sull'intero dataset \emph{prima} della
cross-validation è un errore metodologico noto come \textit{data leakage}:

\begin{itemize}
\item \textbf{Oversampling:} \texttt{Resample} bilancia duplicando istanze della classe
minoritaria. Se applicato prima dello split, le copie vengono distribuite casualmente
nei fold: l'istanza originale X e la sua copia X' --- identiche --- possono finire
rispettivamente nel training fold e nel test fold. Il modello si addestra su X, poi
viene valutato su X', che conosce già esattamente. Le metriche risultano
artificialmente ottimistiche.

\item \textbf{Feature selection:} InfoGain e Spearman assegnano uno score a ogni feature
misurando quanto è correlata con la colonna \texttt{isBuggy}. Se questo calcolo viene
eseguito sull'intero dataset, lo score include anche le istanze che finiranno nel test
fold: si sta barando, perché si usano le etichette di test per decidere quali feature
selezionare, e poi si valuta il modello sulle stesse etichette che hanno già influenzato
quella scelta. Le metriche risultano ottimistiche.
\end{itemize}

La soluzione adottata è \texttt{FilteredClassifier} di Weka, che garantisce che ogni
filtro (balancing e feature selection) venga applicato internamente a ogni fold.
L'ordine dei filtri è rilevante: prima il balancing, poi la feature selection. La FS
deve osservare la distribuzione già bilanciata: calcolare InfoGain su un training
sbilanciato (9\% buggy) produrrebbe score dominati dalla classe maggioritaria.

\paragraph{NONE -- Nessun balancing}

Il classificatore viene addestrato sulla distribuzione originale del training fold
($\approx$9\% buggy, $\approx$91\% non-buggy). Fornisce la baseline di riferimento.

\paragraph{Undersampling -- SpreadSubsample}

Rimuove casualmente istanze dalla classe \textbf{maggioritaria} fino a ottenere un
rapporto 1:1. Implementazione: \texttt{SpreadSubsample(distributionSpread=1.0)}.
Per OpenJPA il training si riduce da $\approx$15.338 a $\approx$2.756 istanze ($-82\%$),
riducendo drasticamente i tempi di addestramento ma eliminando la maggior parte delle
informazioni sulla classe non-buggy.

\paragraph{Oversampling -- Resample}

Ricampiona l'intero dataset con probabilità di estrazione sbilanciate verso una
distribuzione uniforme, senza scartare istanze dalla maggioritaria. Implementazione:
\texttt{Resample(biasToUniformClass=1.0, noReplacement=false)}, con
\texttt{sampleSizePercent} calcolato dinamicamente come
$2 \times N_\text{majority} / N_\text{total} \times 100$. Per OpenJPA il training cresce
da $\approx$15.338 a $\approx$27.000 istanze ($\approx$13.860 per classe). Non aggiunge
nuove informazioni ma ricampiona le istanze minoritarie più frequentemente.

La scelta di \texttt{Resample} rispetto a una duplicazione diretta delle istanze
minoritarie è pragmatica: \texttt{Resample} è l'unico filtro di oversampling supervisionato
disponibile nel jar principale di \texttt{weka-stable 3.8.6} (gli altri filtri supervisionati
di istanza nel core sono \texttt{SpreadSubsample} e \texttt{ClassBalancer}). Duplicare
direttamente le istanze minoritarie richiederebbe un filtro custom, con lo stesso livello
di complessità implementativa affrontato per SMOTE. \texttt{Resample} ottiene il
medesimo risultato pratico --- dataset bilanciato, istanze minoritarie estratte più
frequentemente --- senza codice aggiuntivo.

\paragraph{SMOTE -- Synthetic Minority Over-sampling Technique}

Genera istanze \textbf{sintetiche} della classe minoritaria interpolando tra istanze reali
nello spazio delle feature: per ogni istanza buggy reale, crea nuovi punti lungo il
segmento che la collega ai suoi $k$ vicini più prossimi ($k=5$ di default). SMOTE è
superiore all'oversampling semplice perché non duplica esattamente le stesse istanze,
riducendo il rischio di overfitting e aumentando la variabilità del training set.

\texttt{weka-stable 3.8.6} non include SMOTE nel suo jar principale. Aggiungere il
package ufficiale come dipendenza Maven non è praticabile: non è pubblicato su Maven
Central. Il sorgente ufficiale (\texttt{SMOTE.java}, Ryan Lichtenwalter, Copyright 2008
University of Waikato; algoritmo: Chawla et al.\ 2002, JAIR vol.~16, pp.~321--357) è
stato re-impacchettato nel progetto come \texttt{SmoteFilter}. Le modifiche rispetto
all'originale sono puramente strutturali (package, nome classe, soppressione dei warning
del compilatore); il metodo \texttt{doSMOTE()} che contiene l'intero algoritmo è identico
all'originale.

Per OpenJPA la percentuale di oversampling è $\approx$913\%, portando le istanze buggy
da $\approx$1.378 a $\approx$13.960 (training totale $\approx$27.000 istanze).

\subsubsection{Strategia di validazione: 10-times 10-fold Cross-Validation}

\paragraph{Ruolo dei tre insiemi}

Si distinguono tre insiemi con scopi distinti: il \textbf{training set} (addestrare il
modello), il \textbf{validation set} (selezionare iperparametri --- usato internamente dai
metodi wrapper) e il \textbf{test set} (valutare la performance finale su dati mai visti).
In questo progetto non è presente un test set separato nel senso classico: la 10-fold CV
funge da stima della performance di generalizzazione. In ogni fold, i 9 fold rimanenti
sono il training set e il fold corrente è il test set.

La CV è usata qui nella sua accezione di \textit{stima della performance}, non di
\textit{selezione di iperparametri}: ogni combinazione (classificatore, FS, balancing) è
prefissata in \texttt{config.yaml} e riceve la propria stima indipendente. Il confronto
tra le combinazioni avviene a posteriori, leggendo il CSV di output.

\paragraph{Motivazione della ripetizione a 10 run}

Un singolo run di 10-fold CV produce una stima influenzata dalla divisione casuale
specifica del dataset in fold. Ripetere 10 volte con seed diversi e mediare i risultati
riduce la varianza della stima: le 10 divisioni ``campionano'' lo spazio delle possibili
partizioni, rendendo la stima più stabile e rappresentativa. Il run $i$-esimo usa seed
$i$ sia per lo shuffle che per la \texttt{crossValidateModel} di Weka. Le predizioni
accumulate su tutti i fold e tutti i run vengono usate per calcolare le metriche finali.

\paragraph{Correttezza delle predizioni accumulate}

In 10-fold CV Weka accumula in \texttt{Evaluation} le predizioni di tutti i fold: ogni
istanza viene predetta esattamente una volta da un modello che \emph{non} l'aveva vista
durante il training di quel fold. Accumulare queste $N$ predizioni non costituisce
alcuna forma di leakage: ogni singola predizione è onesta, prodotta su un'istanza che in
quel momento si trovava nel test set.

L'operazione scorretta sarebbe addestrare il modello sull'intero dataset e predire sullo
stesso dataset (training $=$ test): in quel caso il modello conoscerebbe già le risposte.
Nella 10-fold CV ogni istanza è invece predetta da un modello addestrato sui restanti 9
fold, che non l'ha mai vista. Le $N$ predizioni accumulate vengono usate per calcolare le
metriche --- inclusa NPofB20, che richiede l'ordinamento globale di tutte le istanze per
risk score --- in modo più stabile rispetto all'utilizzo delle sole $N/10$ predizioni di
un singolo fold.

\paragraph{FilteredClassifier e prevenzione del data leakage}

\texttt{FilteredClassifier} di Weka garantisce che ogni filtro (balancing e feature
selection) venga applicato \textbf{internamente per ogni fold}, non sull'intero dataset
prima della suddivisione. Il flusso per ogni fold è il seguente:

\begin{quote}
\textit{Fase di training:} \texttt{crossValidateModel} chiama
\texttt{FilteredClassifier.buildClassifier(training\_fold)}. Il filtro di balancing
viene fittato e applicato esclusivamente sul training fold, producendo il training
set bilanciato. La feature selection viene poi fittata sul training fold bilanciato,
calcolando gli score (InfoGain, Spearman) senza visibilità su nessuna istanza del
test fold. Il classificatore base viene addestrato sul risultato finale.
\smallskip

\textit{Fase di predizione:} \texttt{FilteredClassifier.classifyInstance(test\_instance)}
applica esclusivamente la feature selection già fittata (trasformazione, non re-fit).
Il filtro di balancing non viene mai invocato in predizione: i filtri supervisionati
d'istanza (\texttt{Resample}, \texttt{SpreadSubsample}, \texttt{SmoteFilter}) operano
sull'insieme di istanze durante \texttt{buildClassifier} e non hanno semantica su una
singola istanza isolata. Il test fold non subisce né aggiunte sintetiche né rimozioni.
\end{quote}

Non vi è data leakage di alcuna forma: ogni predizione di test è prodotta da un
modello che non ha mai visto quella istanza, con feature selection e balancing calcolati
esclusivamente sui dati di training di quel fold.

\paragraph{Test fold: distribuzione naturale e stratificazione}

Il fatto che ogni fold di test contenga una proporzione di istanze buggy che può
discostarsi da quella del dataset completo non è un problema metodologico: è il
comportamento corretto. Il test set deve riflettere la distribuzione reale del problema.
Il balancing è uno strumento di training --- serve a impedire che il classificatore
apprenda la scorciatoia di predire sempre la classe maggioritaria --- ma non deve
modificare i dati su cui la performance viene misurata. Applicare oversampling o
undersampling anche sul test fold significherebbe valutare su una distribuzione
artificiale, non su quella che il modello incontrerà in produzione: precision e recall
risulterebbero distorti e non interpretabili.

Weka usa \textbf{stratified cross-validation} di default: \texttt{crossValidateModel}
distribuisce le istanze nei fold mantenendo approssimativamente la stessa proporzione
di classi del dataset completo in ogni fold ($\approx$9\% buggy). La stratificazione
riduce la varianza della stima impedendo che un fold di test contenga per caso zero
istanze buggy, il che renderebbe le metriche sulla classe positiva (Precision, Recall,
AUC, NPofB20) indefinite o non confrontabili tra fold.

\paragraph{Cross-validation annidata per i metodi wrapper}
\label{sec:nested-cv}

ForwardSearch e BackwardSearch introducono una struttura di cross-validation annidata: il
loop esterno (10$\times$10 CV) risponde a ``quanto generalizza il modello finale?'',
mentre il loop interno (5-fold CV del wrapper) risponde a ``quali feature rendono
NaiveBayes migliore?''. I due loop rispondono a domande completamente diverse e sono
concettualmente indipendenti. Con \texttt{FilteredClassifier}, il loop interno opera
esclusivamente sui dati del training fold del loop esterno: il test fold del loop esterno
rimane invisibile all'intera procedura di selezione delle feature.

\paragraph{Temporalità: 10$\times$10 CV vs Walk-Forward Validation}
\label{sec:temporal}

La 10$\times$10 CV assegna le istanze ai fold in modo \textbf{casuale}, senza tenere
conto dell'ordine temporale delle release. Un fold di training può quindi contenere
istanze di release più recenti di quelle nel fold di test, introducendo una forma di
\textit{temporal leakage}: il modello apprende pattern che esistono solo perché le classi
sono state modificate o corrette in versioni successive a quelle da predire. Le metriche
risultanti tendono a essere leggermente ottimistiche rispetto all'accuratezza in produzione.

La \textbf{Walk-Forward Validation} elimina questo problema rispettando l'ordine
temporale: ogni iterazione addestra sulle release passate e testa sulla successiva.

\begin{table}[h!]
\centering
\small
\begin{tabular}{@{}p{0.25\textwidth}p{0.33\textwidth}p{0.33\textwidth}@{}}
\toprule
\textbf{Aspetto} & \textbf{10$\times$10 Cross-Validation} & \textbf{Walk-Forward} \\
\midrule
Rispetto della temporalità &
No --- fold casuali, possibile leakage temporale &
Sì --- training sempre nel passato \\
\addlinespace
Numero di valutazioni &
100 (10 run $\times$ 10 fold) &
$\approx$8--9 (una per release) \\
\addlinespace
Varianza della stima &
Bassa --- 100 misure, media stabile &
Alta --- le prime iterazioni hanno pochissimi dati \\
\addlinespace
Correttezza metodologica &
Compromesso accettato in letteratura &
Ideale per simulare l'uso reale \\
\bottomrule
\end{tabular}
\caption{Confronto tra 10$\times$10 Cross-Validation e Walk-Forward Validation.}
\end{table}

Con $\approx$11 release in OpenJPA, la Walk-Forward produrrebbe al massimo 8--9
valutazioni, ciascuna su un test set di una singola release ($\approx$1.000--2.000
istanze). Le prime iterazioni addestrano su 1--2 release sole, rendendo la stima
altamente variabile. La 10$\times$10 CV è il \textbf{compromesso standard} nella
letteratura di bug prediction per dataset con numero limitato di versioni: rinuncia alla
correttezza temporale in favore di una stima più stabile e confrontabile tra
configurazioni diverse. Il temporal leakage è riconosciuto come limitazione e discusso
nella sezione sulle minacce alla validità.

\subsubsection{Metriche di valutazione}

Tutte le metriche sono calcolate sulla \textbf{classe positiva} (buggy = \texttt{true}).
L'accuracy non è adeguata per dataset sbilanciati: un classificatore che predice sempre
``non buggy'' ottiene $\approx$91\% di accuracy senza alcuna utilità pratica. Le cinque
metriche scelte coprono angolazioni complementari e irriducibili: il costo per allarme
(Precision), la completezza (Recall), la qualità del ranking (AUC), l'accordo corretto
per il caso (Kappa) e l'utilità pratica con budget fisso (NPofB20).

\paragraph{Perché Precision e Recall separatamente e non F1}

$F_1 = 2 \cdot (\text{Precision} \cdot \text{Recall}) / (\text{Precision} +
\text{Recall})$ è la media armonica delle due metriche. Collassarle in un unico valore
nasconde però informazioni diagnosticamente rilevanti: se $F_1$ è basso non è possibile
sapere se il problema è una Precision bassa (troppi falsi allarmi), un Recall basso
(troppi bug non trovati) o entrambi. Riportare le due metriche separatamente permette di
vedere esattamente dove il classificatore fallisce e di scegliere il modello in funzione
della priorità del contesto applicativo: se il costo principale sono i bug sfuggiti si
privilegia il Recall, se sono gli effort di ispezione sprecati si privilegia la Precision.
$F_1$ presuppone che le due abbiano uguale importanza, il che non è necessariamente vero
nel dominio della bug prediction.

\paragraph{Precision}

\[
\text{Precision} = \frac{TP}{TP + FP}
\]

Misura il ``costo per allarme'': quando il modello segnala un file come buggy, con quale
probabilità lo è davvero. Precision bassa produce molti falsi allarmi, sprecando effort
di ispezione. Non è sufficiente da sola: alta precision con basso recall significa che il
modello è affidabile nei segnali emessi ma ne mette pochi.

\paragraph{Recall}

\[
\text{Recall} = \frac{TP}{TP + FN}
\]

Misura la ``completezza'': la frazione dei file buggy reali intercettati. Recall bassa
implica che molti bug sfuggono al modello e arrivano in produzione --- il caso più costoso
nello scenario di bug prediction. In un contesto di sicurezza software un falso negativo
(bug non trovato) costa molto più di un falso positivo (file sano ispezionato inutilmente).

\paragraph{AUC -- Area Under the ROC Curve}

Area sotto la curva ROC, che traccia il tasso di veri positivi (Recall) contro il tasso
di falsi positivi ($\text{FPR} = FP/(FP+TN)$) al variare della soglia di classificazione.
$\text{AUC} \in [0,1]$: $0{,}5$ corrisponde al classificatore casuale, $1{,}0$ al
classificatore perfetto. AUC valuta la capacità di \textbf{ordinare} le istanze per
probabilità di bugginess, indipendentemente dalla soglia scelta, ed è invariante allo
sbilanciamento tra classi. Ha anche un'interpretazione probabilistica diretta:
$\text{AUC} = P(\text{score\_buggy} > \text{score\_non-buggy})$, cioè la probabilità
che il modello assegni uno score più alto a un file buggy rispetto a uno sano scelti a caso.

\paragraph{Kappa -- Cohen's Kappa}

Misura di accordo tra predizioni e etichette reali, corretta per l'accordo atteso per
caso. Detti $P_o$ l'accordo osservato e $P_e$ l'accordo atteso per caso:

\[
\kappa = \frac{P_o - P_e}{1 - P_e}
\]

$\kappa = 0$ indica accordo pari al caso; $\kappa = 1$ indica accordo perfetto;
$\kappa > 0{,}6$ è convenzionalmente considerato ``accordo sostanziale'' (Landis \& Koch,
1977). Con $\approx$9\% di buggy, un classificatore sempre-``non buggy'' ha
$P_e \approx 0{,}836$ e $\kappa \approx 0{,}07$: Kappa espone il comportamento degenere
che l'accuracy ($91\%$) maschera completamente.

\paragraph{NPofB20}

Frazione di bug totali che si trovano nel 20\% del codice con risk score più alto:

\[
\text{NPofB20} = \frac{|\text{istanze buggy nel top }20\%|}{|\text{istanze buggy totali}|}
\]

La baseline casuale è $\text{NPofB20} \approx 0{,}20$. Qualsiasi valore superiore indica
che il modello aggiunge valore rispetto al random. NPofB20 traduce la qualità del modello
in termini di effort risparmiato: se un team ispeziona solo il 20\% dei file più a rischio
e $\text{NPofB20} = 0{,}8$, troverà l'80\% di tutti i bug con il 20\% dello sforzo.

NPofB20 non è disponibile in Weka e richiede un calcolo custom: la classe
\texttt{NpofB20Service} estrae le predizioni grezze da \texttt{Evaluation.predictions()},
le ordina per probabilità di bugginess decrescente, e calcola la frazione di istanze buggy
presenti nel primo 20\% dell'ordinamento.

\subsubsection{Configurazione e parallelismo}

\paragraph{Configurazione delle combinazioni}

Le combinazioni da eseguire sono definite in \texttt{config.yaml} come prodotto cartesiano
di tre liste: classificatori, strategie di feature selection e strategie di balancing.
Commentare una voce la esclude dall'esecuzione senza modificare il codice. La
configurazione default esegue 27 combinazioni (3 clf $\times$ 3 FS $\times$ 3 bal, con
SMOTE e metodi wrapper disabilitati per default). Con tutte le strategie abilitate si
arriva a 60 combinazioni.

\paragraph{Parallelismo a due livelli}

La pipeline supporta due livelli indipendenti di parallelismo, configurabili in
\texttt{config.yaml}:

\begin{enumerate}
\item \textbf{Livello combinazione:} ogni combinazione è completamente indipendente dalle
altre (nessuno stato condiviso modificabile tra thread). Le combinazioni vengono
distribuite su un thread pool fisso e i risultati raccolti tramite \texttt{Future}.

\item \textbf{Livello interno RandomForest:} Weka's \texttt{RandomForest} supporta la
costruzione parallela degli alberi tramite \texttt{setNumExecutionSlots(k)}. Non influisce
su NaiveBayes e IBk.
\end{enumerate}

I due livelli si \textbf{moltiplicano}: il numero di core usati è approssimativamente
$\texttt{combinations} \times \texttt{randomForestSlots}$. I parametri accettano valori
numerici espliciti (\texttt{"4"}), percentuali (\texttt{"50\%"}) o \texttt{"auto"} (tutti
i core disponibili). È disponibile anche una modalità interattiva che mostra le risorse
disponibili e chiede i valori a runtime.

\begin{table}[h!]
\centering
\small
\begin{tabular}{@{}lrr@{}}
\toprule
\textbf{Configurazione} & \textbf{Combinazioni} & \textbf{Stima (4 thread)} \\
\midrule
Default (3 clf $\times$ 3 FS $\times$ 3 bal, SMOTE disabilitato) & 27 & $\approx$35 min \\
+ SMOTE (3 clf $\times$ 3 FS $\times$ 4 bal) & 36 & $\approx$1 ora \\
+ ForwardSearch o BackwardSearch & $+$9 ciascuno & $+$9--27 ore \\
Tutto abilitato (3 clf $\times$ 5 FS $\times$ 4 bal) & 60 & 1--3 giorni \\
\bottomrule
\end{tabular}
\caption{Stime di runtime calibrate su OpenJPA ($\approx$15.338 istanze, rfSlots$=1$) su
hardware con CPU mid-range. I tempi base per classificatore (RF $\approx$757\,s, IBk
$\approx$227\,s, NB $\approx$6\,s) sono hardware-specifici; i moltiplicatori di balancing
(Undersampling $\approx$0{,}20$\times$, Oversampling $\approx$1{,}65$\times$) sono più
stabili tra hardware diversi.}
\end{table}

\subsubsection{Note implementative rilevanti}

\paragraph{Correzione del double-randomization bug}

\texttt{Evaluation.crossValidateModel(classifier, data, folds, random)} gestisce
internamente shuffle e stratificazione prima di costruire i fold. Una versione precedente
della pipeline pre-shufflava il dataset con \texttt{data.randomize(new Random(run))}
prima di passarlo a \texttt{crossValidateModel}: il metodo Weka applicava poi il proprio
shuffle con lo stesso seed, producendo fold diversi da quelli intesi (double-randomization)
e alterando silenziosamente la riproducibilità. La correzione consiste nel passare il
dataset originale direttamente a \texttt{crossValidateModel}, lasciando a Weka l'unica
responsabilità della randomizzazione.

\paragraph{Componenti custom: SmoteFilter e SpearmanAttributeEval}

Due componenti del package \texttt{com.isw2project.classifier} non erano disponibili nel
jar principale di \texttt{weka-stable 3.8.6}:

\begin{itemize}
\item \texttt{SmoteFilter}: re-package del sorgente ufficiale Weka SMOTE (Ryan
Lichtenwalter, 2008, University of Waikato). Il package ufficiale non è pubblicato su
Maven Central e non può essere dichiarato come dipendenza Maven standard. Le modifiche
rispetto all'originale sono puramente strutturali; il metodo \texttt{doSMOTE()}, che
contiene l'intero algoritmo di interpolazione, è identico all'originale.

\item \texttt{SpearmanAttributeEval}: implementazione originale scritta da zero.
\texttt{weka-stable} include \texttt{CorrelationAttributeEval} (Pearson) ma non un
evaluator per Spearman. Pearson assume una relazione lineare ed è sensibile agli outlier:
entrambi problemi critici per metriche di codice con distribuzioni fortemente skewed come
\texttt{LOC\_Touched} o \texttt{Churn}.
\end{itemize}

\paragraph{Gestione dei valori NaN nelle metriche}

\texttt{Evaluation.precision(classIndex)} e metodi analoghi restituiscono
\texttt{Double.NaN} quando una classe non produce predizioni positive in nessun fold
(tipicamente quando il classificatore non riesce mai a predire la classe minoritaria).
Un \texttt{NaN} si propaga aritmeticamente attraverso tutte le somme successive,
rendendo l'intera riga del CSV inutilizzabile. Il metodo \texttt{safe(double v)} in
\texttt{ClassifierEvaluatorService} intercetta \texttt{NaN} e valori infiniti
sostituendoli con \texttt{0.0} prima dell'accumulo, garantendo che ogni combinazione
produca comunque una riga leggibile nel CSV.

\paragraph{Configurazione di \texttt{WrapperSubsetEval}: ordine delle chiamate API}

\texttt{WrapperSubsetEval} implementa l'interfaccia \texttt{OptionHandler} di Weka,
che prescrive che \texttt{setOptions()} ripristini \emph{tutti} i parametri ai valori
di default prima di applicare le opzioni ricevute. Chiamare
\texttt{setClassifier(NaiveBayes)} \emph{prima} di \texttt{setOptions(\{"-E","AUC"\})}
produce quindi un effetto inatteso: internamente, \texttt{setOptions} azzera il
classificatore al suo default (\texttt{ZeroR}) prima ancora di processare le opzioni
passate. Il risultato è che il classificatore interno effettivo diventa \texttt{ZeroR},
non \texttt{NaiveBayes}.

Con \texttt{ZeroR} come classificatore interno, ogni sottoinsieme di feature produce
la stessa performance: \texttt{ZeroR} predice sempre la classe maggioritaria
indipendentemente dalle feature presenti, ottenendo AUC$=0{,}5$ su qualsiasi subset.
\texttt{GreedyStepwise} con ricerca backward parte da tutte le 19 feature e rimuove
iterativamente quella la cui eliminazione non abbassa la metrica interna; poiché la
metrica è costante al variare del sottoinsieme, il processo converge al sottoinsieme
vuoto. Il risultato osservato era il messaggio Weka
\textit{``Cannot build model (only class attribute present in data!)''}
per tutte le combinazioni BackwardSearch.

La soluzione consiste nell'invertire l'ordine delle chiamate: \texttt{setOptions()}
va chiamato \emph{prima} di \texttt{setClassifier()}, in modo che \texttt{NaiveBayes}
venga impostato dopo il reset interno e sopravviva. Con \texttt{NaiveBayes} come
classificatore interno, sottoinsiemi di feature diversi producono performance
misurabili e differenti, e la ricerca greedy converge a un sottoinsieme non vuoto.
Questa è una proprietà documentata del contratto \texttt{OptionHandler} di Weka,
non un bug della libreria.

Come metrica interna è stato adottato \textbf{AUC} (opzione \texttt{-E AUC}) invece
dell'accuracy di default. Su un dataset sbilanciato come OpenJPA
($\approx$91\% non-buggy), l'accuracy è dominata dalla classe maggioritaria e meno
sensibile nel discriminare subset migliori da subset peggiori; AUC misura la capacità
discriminante indipendentemente dalla soglia di decisione ed è più indicativo su
distribuzioni sbilanciate. La scelta è coerente con la metrica di valutazione primaria
adottata per i classificatori in questo studio.

\paragraph{Soppressione dei warning}

Il progetto sopprime quattro categorie di warning, ognuna per una ragione distinta:

\begin{itemize}
\item \textbf{JVM-level --- \texttt{.mvn/jvm.config}: \texttt{--enable-native-access=ALL-UNNAMED}.}
Da Java~24 la JVM restringe l'accesso ai metodi nativi ai soli moduli esplicitamente
autorizzati. Weka usa \texttt{jniloader} per caricare librerie native di algebra lineare
(\texttt{ARPACK}, \texttt{BLAS}); poiché \texttt{jniloader} è in un modulo unnamed,
senza il flag la JVM emette un warning ad ogni esecuzione. Il flag è in
\texttt{.mvn/jvm.config} anziché in \texttt{pom.xml} per non inquinare il build
descriptor con configurazione JVM runtime-specifica.

\item \textbf{Codice legacy --- \texttt{SmoteFilter.java}:
\texttt{@SuppressWarnings(\{"unchecked","rawtypes","deprecation"\})}.}
Il sorgente di \texttt{SmoteFilter} è il codice originale Weka del 2008, scritto in
stile Java pre-generics. I tre warning sono attesi e inevitabili senza riscrivere
l'algoritmo: \texttt{rawtypes} per le collezioni senza tipo generico, \texttt{unchecked}
per i cast da raw type, \texttt{deprecation} per \texttt{new Double(String)} rimosso in
Java~9. La soppressione garantisce che il metodo \texttt{doSMOTE()} rimanga identico
all'originale.

\item \textbf{SonarQube S1172 --- \texttt{Milestone1Main} e \texttt{Milestone2Main}:
\texttt{@SuppressWarnings("java:S1172")}.}
La regola segnala \texttt{String[] args} di \texttt{main} come parametro inutilizzato.
La firma \texttt{public static void main(String[] args)} è imposta dalla JVM come entry
point e non può essere modificata indipendentemente dall'uso di \texttt{args}.

\item \textbf{SonarQube S106 --- \texttt{InteractiveParallelismConfigurator}:
\texttt{@SuppressWarnings("java:S106")}.}
La regola vieta l'uso diretto di \texttt{System.out}, raccomandando un logger. Questa
classe mostra un menu interattivo e legge la risposta da stdin: usare un logger
introdurrebbe timestamp e livelli di log, rompendo l'esperienza interattiva.
\texttt{System.out} è la scelta corretta; la soppressione documenta l'intenzionalità.
\end{itemize}

\paragraph{Code smell di \texttt{SmoteFilter}: scelta consapevole di non intervenire.}
L'analisi statica con SonarQube rileva 29 code smell in \texttt{SmoteFilter.java}.
Questi non sono stati corretti per una ragione precisa: \texttt{SmoteFilter} non è
codice scritto nel contesto di questo progetto, ma il re-package del sorgente ufficiale
Weka SMOTE (Ryan Lichtenwalter, 2008), incluso nel progetto esclusivamente perché non
disponibile su Maven Central. Le uniche modifiche apportate rispetto all'originale sono
puramente strutturali: cambio del package, rinomina della classe e aggiunta dell'annotazione
\texttt{@SuppressWarnings(\{"unchecked","rawtypes","deprecation"\})} per sopprimere i
warning del compilatore dovuti allo stile Java pre-generics del 2008.

Correggere i code smell di \texttt{SmoteFilter} significherebbe riscrivere codice esterno
non di proprietà, con il rischio di introdurre regressioni nell'algoritmo. La scelta
coerente è trattarla come una dipendenza di terze parti: ci si affida all'originale
verificato e si limitano le modifiche al minimo strutturale necessario per l'integrazione.
